\documentclass[11pt]{article}
\usepackage[a4paper,margin=24mm]{geometry}
\usepackage{amsmath,amssymb,amsthm,booktabs,hyperref}
\newtheorem{theorem}{Theorem}\newtheorem{lemma}[theorem]{Lemma}
\newtheorem{corollary}[theorem]{Corollary}
\newcommand{\N}{\operatorname{N}}\newcommand{\Inn}{\operatorname{Inn}}
\title{Transport, incidence weights, and certificates for finite subgroup data}
\author{A component for the subgroup atlas of Connor and Leemans}
\date{}
\begin{document}\maketitle
\noindent
We prove the correctness of the two algorithmic constructions in Sections 2
and 3 of \cite{CL}, using its seven-page author edition dated 2 December
2013. The maximal-subgroup descent and the recursive choice of normal
subgroups are the authors' ideas. Our contribution here is an explicit
certificate interface: changes of permutation action, fusion of conjugacy
classes, incidence weights, and extension data have separate hypotheses
and conclusions. We do not claim that these elementary arguments are new
group-theoretic results. In particular, this component does not reproduce
the more than one hundred numerical lattices announced in Section 4.

\paragraph{Conventions and inputs.}
All groups are finite. Write $H^g=g^{-1}Hg$ and $[H]_G$ for the conjugacy
class of a subgroup. Order the classes by $[K]_G\leq[H]_G$ when some
conjugate of $K$ is contained in $H$. This is a partial order: a relation
in both directions forces equal orders, and containment between finite
subgroups of equal order is equality. We use the term \emph{poset}; none
of the arguments assumes that this poset is a lattice. In the actual
subgroup lattice the join is $\langle H,K\rangle$, not set union.

For each subgroup $H$ selected for processing, the input consists of an
isomorphism $\phi_H:H\longrightarrow\widetilde H$ and a complete list of
representatives of the $\widetilde H$-conjugacy classes of maximal proper
subgroups. Completeness of this list is an explicit input, not a consequence
of the orders or number of the returned subgroups. Conjugacy tests in the
ambient group $G$ must be exact. A positive conjugacy test can be checked
by its transporter; a negative test requires a complete test or certificate.

\begin{lemma}\label{transport}
An isomorphism $\phi:H\to\widetilde H$ transports subgroup containment,
maximality, conjugacy within $H$, and normalizers. In particular, if
$\widetilde M$ runs through a complete set of maximal-subgroup class
representatives, then $\phi^{-1}(\widetilde M)$ does so in $H$, with
\[
 [H:\N_H(\phi^{-1}(\widetilde M))]
   =[\widetilde H:\N_{\widetilde H}(\widetilde M)].
\]
\end{lemma}
\begin{proof}
The mutually inverse maps on subgroups are image and preimage. They preserve
containment in both directions, hence preserve the absence of a strictly
intermediate subgroup. The identity $\phi(M^h)=\phi(M)^{\phi(h)}$ transports
conjugacy. Applying it with $M^h=M$ identifies the two normalizers; their
indices agree because both group orders agree.
\end{proof}

For a reduced permutation action, it is therefore necessary to establish
faithfulness, rather than merely the existence of an action homomorphism.
An explicit inverse is sufficient. So is a presentation of its image on
the actual images of the generators, provided that all defining relators
hold on the original generators: the resulting reverse homomorphism is
inverse on generators and consequently on the groups. This last criterion
is proved and checked in the separate faithful-action component of this
project; it is an optional way of supplying the isomorphism required here.

\begin{theorem}\label{descent}
Start with a registry containing $[G]_G$ and a queue containing $G$.
Process each registered class once. For a representative $H$, transport
its complete maximal-subgroup list by Lemma~\ref{transport}, identify each
resulting subgroup up to $G$-conjugacy, and register and queue every class
not already registered. Record an edge from $[H]_G$ to each resulting
class, retaining separately all local $H$-conjugacy classes that fuse there.
The process terminates and returns every subgroup class of $G$ exactly
once. Reachability in the resulting graph is precisely containment in
the class poset.
\end{theorem}
\begin{proof}
There are finitely many subgroups, and a registered class is never queued
again. Thus termination holds even if the queue is processed in arbitrary
order. Every inserted vertex and edge is an actual subgroup class and an
actual proper containment.

Given $K\leq H$, refine the interval from $K$ to $H$ to a chain of maximal
proper inclusions; this is possible by finiteness. Induction down such a
chain proves that a conjugate of $K$ is discovered whenever the class of
$H$ is processed. At each step the complete local maximal list contains a
conjugate within the current parent, and fusion changes that representative
only by conjugacy in $G$. Starting from $G$ proves completeness.

Edges imply containment and containment is transitive on classes.
Conversely, apply the same maximal-chain argument inside a representative
of the larger class to obtain a path to the smaller class. Exact ambient
conjugacy tests and the registry ensure that vertices are neither duplicated
nor incorrectly identified.
\end{proof}

The registry in this formulation contains both queued and processed
classes. Testing only the current queue, as the brief prose algorithm in
\cite{CL} suggests literally, needs an additional scheduling invariant to
avoid rediscovering a processed class. The permanent registry makes that
invariant unnecessary. A computation begun from partial data must still
certify that every registered parent has a complete processed maximal list;
unprocessed vertices remain obligations.

\paragraph{Weights after fusion.}
Fix representatives $H,K$. Let $M_1,\ldots,M_t$ be the representatives
of the $H$-conjugacy classes of maximal subgroups of $H$ which are
$G$-conjugate to $K$. Put
\[
 a^{\max}_{H,K}=\sum_{j=1}^t[H:\N_H(M_j)],\qquad
 c_H=[G:\N_G(H)].                                      \tag{1}
\]
An empty sum is zero. Several local classes may fuse into one ambient
class; replacing this sum by one summand is not valid.

\begin{theorem}\label{weights}
The number $a^{\max}_{H,K}$ counts the maximal subgroups of the actual
$H$ belonging to $[K]_G$. The number of conjugates of $H$ in which an
actual $K$ is maximal is
\[
 b^{\max}_{K,H}=\frac{c_Ha^{\max}_{H,K}}{c_K}.
                                                               \tag{2}
\]
The cover relations of the class poset are the edges left after transitive
reduction of the graph in Theorem~\ref{descent}. For such a cover,
$a^{\max}_{H,K}$ and $b^{\max}_{K,H}$ are the full incidence numbers
$n_{HK}$ and $n_{KH}$ defined in Section 2 of \cite{CL}.
\end{theorem}
\begin{proof}
Orbit--stabilizer for the conjugation action of $H$ gives exactly
$[H:\N_H(M_j)]$ distinct subgroups in the $j$th local class. These local
classes are disjoint and complete, proving (1). Count the pairs
$(H',K')$ with $H'\in[H]_G$, $K'\in[K]_G$, and $K'$ maximal in $H'$.
Counting by the first coordinate gives $c_Ha^{\max}_{H,K}$; counting by
the second gives $c_Kb^{\max}_{K,H}$. Conjugacy makes each fibre size
constant, proving (2), including its integrality.

The graph is acyclic because orders strictly decrease along edges. Its
reachability is the class order, so its transitive reduction gives exactly
the covers. If $[H]_G$ covers $[K]_G$ and $K'\leq H'$ are representatives,
any strict intermediate subgroup would give a strict intermediate class.
Thus every such containment is maximal and the maximal-incidence counts
are the required full incidence counts.
\end{proof}

For a pair that is not a cover, (1) only counts maximal embeddings.
It need not count every embedding of the smaller class into the larger.
Likewise, normalizer orders alone do not replace the local fusion data.
The preceding results recover the weighted class poset without asserting
that it determines $G$ up to isomorphism.

For example, take $G=S_4$, $H=\langle(1\,2),(3\,4)\rangle$, and
$K=\langle(1\,2)\rangle$. The two subgroups generated by the individual
transpositions are distinct local classes in the abelian group $H$, but
they are conjugate in $G$. Hence $a^{\max}_{H,K}=2$, $c_H=3$, $c_K=6$,
and $b^{\max}_{K,H}=1$. Keeping only one of the two local classes would
give the wrong weight. This example occurs in the literal $S_4$ check.

\paragraph{What a recursive group name certifies.}
Let $N\lhd H$, let $Q=H/N$, and choose a section $s:Q\to H$ with
$s(1)=1$. Define
\[
 \alpha_q(n)=s(q)ns(q)^{-1},\qquad
 f(q,r)=s(q)s(r)s(qr)^{-1}\in N.                         \tag{3}
\]
Here $\Inn(u)(n)=unu^{-1}$, and the multiplication on $N$ and $Q$ is part
of the data. A structural
expression involving only their names and a dot omits the maps in (3).

\begin{theorem}\label{extension}
The map $(n,q)\mapsto ns(q)$ is a bijection $N\times Q\to H$, and under
this bijection the multiplication is
\[
 (n,q)(m,r)=(n\alpha_q(m)f(q,r),qr).                     \tag{4}
\]
Conversely, automorphisms $\alpha_q$ and elements $f(q,r)$ define a group
by (4), with kernel $N$ and quotient $Q$, if
\begin{align*}
 \alpha_1&=\mathrm{id},&f(1,q)&=f(q,1)=1,\\
 \alpha_q\alpha_r&=\Inn(f(q,r))\alpha_{qr},\qquad&
 f(q,r)f(qr,t)&=\alpha_q(f(r,t))f(q,rt).
\end{align*}
Thus the displayed data, with exact identifications of kernel and quotient,
certify the extension, not merely its order or composition factors.
\end{theorem}
\begin{proof}
Every $h\in H$ lies in exactly one quotient coset $q$ and has a unique
expression $ns(q)$; this proves the bijection. Inserting $s(q)^{-1}s(q)$
in a product gives (4). Conjugation by $s(q)s(r)=f(q,r)s(qr)$ proves the
first identity involving $\alpha$. Associativity of $s(q)s(r)s(t)$ proves
the last identity.

Conversely, compute the two associations of three pairs in (4). The
identity involving $\alpha$ moves the conjugating factor $f(q,r)$ through
$\alpha_{qr}$; the last identity then makes the two products equal.
The normalization gives identity $(1,1)$. Given $(n,q)$, its product with
$(m,q^{-1})$ is the identity for the unique
$m=\alpha_q^{-1}(n^{-1}f(q,q^{-1})^{-1})$. An associative finite monoid
in which each element has a right inverse is a group: if $xy=1$ and
$yz=1$, then $x=x(yz)=(xy)z=z$, hence $yx=1$. Finally the projection to
$Q$ is a surjective homomorphism and its kernel is $N\times\{1\}$.
\end{proof}

An explicit complement $C$ with $N\cap C=1$ and $NC=H$ certifies a split
extension; its inverse quotient map supplies a section with $f=1$.
A complement centralizing $N$ certifies a direct product. Conversely, a
homomorphic section supplies a complement, and a homomorphic section whose
image centralizes $N$ supplies a direct decomposition. Failure to find a
complement does not certify a nonsplit extension: that conclusion needs
an exhaustive complement test or a separate obstruction proof.

For instance, let $H=\langle x\rangle\cong C_4$ and $N=\langle x^2\rangle$.
The normalized section taking the nonidentity quotient element to $x$
centralizes $N$, but its factor at that pair is $x^2\ne1$. It is not a
homomorphic section. Replacing this factor by $1$ gives $C_2\times C_2$;
the kernel and quotient have not changed, but the multiplication has.
The finite check rejects that replacement as a certificate for the original
$C_4$, even though it is valid data for a different extension.

\begin{corollary}
The recursive naming procedure of Section 3 of \cite{CL} terminates if,
at each nonterminal step, it chooses a proper nontrivial normal subgroup
and recursively treats its kernel and quotient. Exact identifications at
the leaves and exact extension data at the internal vertices determine the
original group. Omitting the action and factor data gives only a structural
description, and distinct groups can have the same resulting description.
\end{corollary}
\begin{proof}
Both child orders are smaller than the parent order, so induction on order
proves termination. At a terminal simple group or a selected recognizable
group the supplied identification is checked directly. At each internal
vertex Theorem~\ref{extension} reconstructs the parent from the already
identified children and the recorded multiplication. Induction proves the
claim for the root. No uniqueness claim is made for descriptions that
discard this multiplication data.
\end{proof}

\paragraph{Evidence and limits.}
The accompanying finite checks compare maximal descent with an independent
enumeration obtained by adjoining elements to subgroups. They check local
fusion, both incidence counts, transitive reduction, and extension
multiplication on small explicit groups. These checks validate the stated
interfaces and supply examples. They do not establish completeness of
the external maximal-subgroup algorithms on every atlas group, reproduce
the atlas's numerical files, or reproduce its historical timing comparison.
The full L054 manuscript therefore remains pending. No conjecture,
compression ratio, or Lean verification is claimed by this component.

\begin{thebibliography}{1}
\bibitem{CL} T. Connor and D. Leemans,
\emph{An Atlas of subgroup lattices of finite almost simple groups},
arXiv:1306.4820v2, 29 November 2013; retained author PDF dated
2 December 2013, Sections 2--4.
\end{thebibliography}
\end{document}
